\chapter{Environmental and Sustainability Analysis}
\label{sec:environmental}

An instrument built for environmental monitoring should itself leave a small
footprint, so this chapter considers both the problem the OpenStat addresses
and its impact over its life.

Antiretrovirals such as efavirenz pass largely untreated through South
African wastewater treatment works and reach surface water at up to
\SI{102}{\micro\gram\per\litre}~\cite{vogt_arv_sa_waters}, far above the
\si{\nano\gram\per\litre} levels shown to harm fish~\cite{env_robson2017}.
Monitoring is limited by the cost of laboratory analysis, which a low-cost
portable instrument helps to lower. The components for the prototype cost
about R2\,230 including VAT (Appendix~\ref{app:bom}).

The largest savings were made in the hardware itself. The AD5941 integrates all the elements discussed in Section \ref{sec:potentiostat_fundamentals} into one
millimeter-scale package~\cite{ad5941_ds},
keeping the board small and the part count low. The SPEs, the only
consumables, were cleaned using ethanol and reused throughout
testing. Table~\ref{tab:env_lifecycle} lists the remaining measures.

\begin{table}[htbp]
  \caption{Further life-cycle measures.}
  \label{tab:env_lifecycle}
  \centering
  \small
  \begin{tabularx}{\textwidth}{@{}l >{\raggedright\arraybackslash}X@{}}
    \toprule
    \textbf{Stage} & \textbf{Design response} \\
    \midrule
    Energy       & \SI{0.89}{\watt} measured load
                   (Eq.~\eqref{eq:hw_pbat}) which equates to $\approx 1500$ scans per charge \\
    Obsolescence & Web interface and display which are open-source and can be updated \\
    End of life  & Recycling under the EPR Regulations~\cite{env_epr2020} \\
    \bottomrule
  \end{tabularx}
\end{table}

A full charge draws about \SI{29}{\watt\hour} (including charging losses) from a grid emitting
\SI{0.960}{\kilo\gram}~CO$_2$e/\si{\kilo\watt\hour}~\cite{env_dffe_gef2024},
roughly \SI{28}{\gram} CO$_2$e per charge or \SI{0.02}{\gram} per scan, so
operating emissions are negligible.
% CHECK: 29 Wh includes the linear TP4056 (~3.7 V / 5 V) and an ~85 % USB
%        adapter.

The main remaining improvement is in the power chain. Boosting the battery to
\SI{5}{\volt} and then dropping to \SI{3.3}{\volt} discards about a third of
the stored energy. By using a \SI{3.3}{\volt} regulator the run time would extend by roughly
\SI{50}{\percent} with the same cell. Safety, cost and societal
considerations are discussed in Appendix~\ref{app:ga3}.


% CHECK: 50 % = (5 V x 160 mA / 0.9) / (3.3 V x 160 mA / 0.9) = 1.52.

%5 Ah × 3.7 V = 18.5 Wh

%2. Loss in the TP4056 charger

%The TP4056 is a linear charger. It takes current from the 5 V USB input and passes %the same current into the battery, burning off the voltage difference as heat. Its %efficiency is therefore roughly the ratio of the two voltages:

%η_charger ≈ 3.7 V / 5 V = 0.74

%3. Loss in the USB adapter

%A typical wall adapter is about 85 % efficient. This is the assumption flagged with % CHECK in the file.

%4. Combine the losses

%η_total = 0.74 × 0.85 ≈ 0.63

%Energy from the wall = 18.5 Wh / 0.63 ≈ 29 Wh